\documentclass[11pt]{article}
\usepackage[a4paper,margin=25mm]{geometry}
\usepackage{amsmath,amssymb,amsthm,mathtools}
\usepackage[T1]{fontenc}
\usepackage{lmodern}
\usepackage{hyperref}
\usepackage{enumitem}
\newtheorem{theorem}{Theorem}[section]
\newtheorem{lemma}[theorem]{Lemma}
\newtheorem{proposition}[theorem]{Proposition}
\newtheorem{corollary}[theorem]{Corollary}
\theoremstyle{definition}
\newtheorem{definition}[theorem]{Definition}
\newtheorem{example}[theorem]{Example}
\newcommand{\card}[1]{\lvert #1\rvert}
\newcommand{\CL}{C_{\mathrm L}}
\newcommand{\CM}{C_{\mathrm M}}
\newcommand{\R}{\mathcal R}
\newcommand{\N}{\mathcal N}
\title{Resource-indexed Candidate Coverage for Conditional Planning\\
\large Consolidated definitions, proofs, separation results, and limitations}
\author{SCAR research working supplement}
\date{16 October 2026}
\begin{document}
\maketitle
\begin{abstract}
We formalize the successful-candidate coverage bounded by the implemented
scalar, lineage, and matching recurrences. The supplement supplies a hereditary
soundness proof, a separation from the conjunction of complete proper-subset
tests and an explicit resource-aware information control, a Hall-deficiency
characterization of matching rejection, safe partial-evaluation rules, and
search and complexity guarantees. A concrete planning counterexample disproves
unconditional strengthening under added projection predicates. A restricted
simulation theorem and a conservative envelope replace that unsupported claim.
These are mathematical results under explicit assumptions, not evidence of
runtime improvement or priority over every labelled planning abstraction.
\end{abstract}

\section{Problem, implementation scope, and the quantity being bounded}
Let $I$ be a finite set of $N$ named initial candidates. A name denotes one
initial state, not a stochastic sample or a posterior probability. Determinism
gives at most one current state $x_m$ per surviving name $m$ along an observation
history. A state consists of Boolean facts and remaining resources
$r\in\mathbb N^b$. Facts immutable under every action are fixed by the name.
For every grounded action $a$, its Boolean precondition, deterministic simultaneous
effects $T_a$, optional Boolean observation $O_a$ evaluated \emph{after} effects,
and fixed consumption $c(a)\in\mathbb N^b$ are given. Applying an action requires
its precondition and $r\geq c(a)$ componentwise; resources become $r-c(a)$.
Actions have unit length. Resources occur only in affordability: there is no
replenishment, numeric goal, or resource-dependent Boolean effect/observation.
The legality convention on conflicting Boolean effects must be the same in
the transition interpreter and channel formula (the current interpreter gives
adds precedence over deletes).

A belief $B=\{(m,x_m):m\in M\}$ is exact. A strong policy chooses one legal
action for every state in the current belief, branches on \emph{actual}
observations, and stops only when every surviving state satisfies the goal.
Its cost is its largest root-to-terminal action count. Write $V(B)$ for the
minimum such cost, or $\infty$ if there is no finite strong policy. The bounded
decision problem is $V(B)\leq d$. Goal truth is not implicitly observed;
already-goal states remain in the exact belief and can constrain subsequent
actions. A proof of $V(B)>H$ is not a proof of $V(B)=\infty$.

Fix a set $P$ of retained dynamic atoms. Put $q=(r,z)$ where $z=x|_P$.
The current implementation applies a coverage query only when all
\emph{unfinished} worlds share $q$. Otherwise it skips the query. For name $m$,
let $\mathcal X_m(q)$ be all non-goal Boolean states with its immutable facts,
retained values $z$, and resources $r$. Forgotten dynamic facts are unrestricted;
reachability and hidden invariants are not assumed. Thus this set includes the
actual state whenever the query is applicable.

\begin{definition}[Hereditary existential coverage]
For $S\subseteq I$, define $\Gamma(q,S,d)$ as the maximum $\card T$ over subsets
$T\subseteq S$ for which states $x_m\in\mathcal X_m(q)$, $m\in T$, and a
depth-at-most-$d$ strong policy for those states exist. The empty subset is
allowed and has value zero. State witnesses may be chosen separately per name.
\end{definition}
This definition is deliberately more permissive than the actual belief. It
also avoids a gap in a proof stated only for policies that succeed on \emph{all}
names of a recursive mask: a child mask can strictly contain the names actually
entering that child. The proof below bounds every such successful subset.

\section{Channels and three complete recurrences}
For an affordable $a$, define the complete existential channel table
\begin{align*}
G_a(q)&=\{m:\exists x\in\mathcal X_m(q),\ a\text{ legal at }x,
                                      \ T_a(x)\models g\},\\
E_{aj}(q)&=\{m:\exists x\in\mathcal X_m(q),\ a\text{ legal at }x,
               \ T_a(x)\not\models g,\ (O_a(x),T_a(x)|_P)=j\}.
\end{align*}
Here $j=(o,z')$ includes the post-effect observation; for a non-observing action
$o$ is a single dummy symbol. Its successor key is $q_{aj}=(r-c(a),z')$.
An individual name can belong to several $E_{aj}$, or both $G_a$ and an $E_{aj}$,
because different forgotten-state witnesses can establish membership.
Empty buckets can be removed. The bound grants knowledge of $z'$ in addition
to $o$, whereas a real policy does not obtain that knowledge for free.

Set all depth-zero capacities to zero, and use zero for an empty action maximum.
The global scalar recurrence is
\begin{equation}\label{eq:scalar}
C_0(q,d)=\min\left(N,\max_a\left[\card{G_a}+
       \sum_j\min\{\card{E_{aj}},C_0(q_{aj},d-1)\}\right]\right).
\end{equation}
For a mask $S$, put $F_a=S\cap G_a$, $R_a=S\setminus G_a$, and
$S_{aj}=R_a\cap E_{aj}$. Define
\begin{equation}\label{eq:lineage}
\CL(q,S,d)=\min\left(\card S,\max_a\left[
                   \card{F_a}+\sum_j\CL(q_{aj},S_{aj},d-1)\right]\right).
\end{equation}
The immediate absorption of \emph{all} $F_a$ is optimistic, even if their actual
states would not finish. Once counted, their names do not recur below that
action. This removes one source of repeated counting.

For matching, form a bipartite graph with unit-capacity model vertices $R_a$,
bucket vertices $j$, edges $m\to j$ when $m\in S_{aj}$, and bucket capacities
$c_j=\CM(q_{aj},S_{aj},d-1)$. Let $\nu_a(c)$ be its maximum integral matching
cardinality with these bucket capacities. Define
\begin{equation}\label{eq:matching}
\CM(q,S,d)=\min\left(\card S,\max_a[\card{F_a}+\nu_a(c)]\right).
\end{equation}
Both relational capacities are zero for $S=\varnothing$. The older implementation
also takes a minimum with $C_0$; the next result proves that ceiling redundant
for complete, unlimited recurrences. A budget fallback is a different evaluation
and must not be used to infer equality of partial execution traces.

\begin{lemma}[Ordering and the redundant ceiling]\label{lem:order}
At the same channels and depth,
$0\leq\CM(q,S,d)\leq\CL(q,S,d)\leq\min\{\card S,C_0(q,d)\}$.
\end{lemma}
\begin{proof}
Induct on $d$, with zero base values. By induction each lineage child is at
most $\min\{\card{S_{aj}},C_0(q_{aj},d-1)\}$, hence at most the corresponding
scalar child contribution. Also $\card{F_a}\leq\card{G_a}$ and $\card S\leq N$.
Taking the action maximum and cap proves the lineage inequality. Matching is
monotone in bucket capacities, since increasing a capacity preserves every
feasible allocation. Its value is at most the sum of bucket capacities.
Inductively replacing matching child values by lineage child values therefore
bounds each matching action by its lineage action. The cap preserves the result.
\end{proof}

\begin{lemma}[Monotonicity in the quantities that are actually ordered]
For either relational recurrence, $S\subseteq S'$ implies
$C(q,S,d)\leq C(q,S',d)$. Capacities are nondecreasing in $d$ and componentwise
in $r$ when all other data are fixed. This statement does not order projections.
\end{lemma}
\begin{proof}
Induct on depth. Enlarging $S$ enlarges $F_a$, $R_a$, and each child mask;
child capacities do not decrease by induction. An old matching embeds into
the enlarged network. Both the action maximum and its cardinality cap are
nondecreasing. For depth, use the zero base and apply induction to each child.
For resources, the templates $G_a,E_{aj},z'$ do not depend on $r$; larger $r$
permits every old action and leaves larger child resources. Induction again
orders child values and action contributions. Scalar monotonicity follows in
the same way. No claim of marginal subadditivity in the mask is required.
\end{proof}

\begin{theorem}[Hereditary soundness]\label{thm:sound}
$\Gamma(q,S,d)\leq\CM(q,S,d)\leq\CL(q,S,d)\leq C_0(q,d)$.
\end{theorem}
\begin{proof}
The assertion is trivial at depth zero because all chosen states are unfinished.
Assume it at $d-1$. Take any successful subset $T\subseteq S$, its state
witnesses, and a policy of depth at most $d$. Its first action $a$ is affordable.
Credit every name in $F_a$, even those absent from $T$. For each successful
name in $T\setminus G_a$, its actual successor is non-goal; otherwise it would
belong to $G_a$. Its successor and observation establish membership in exactly
one actual bucket $j$ and in $S_{aj}$.

Let $T_j$ be these successful names in bucket $j$. The actual continuation,
restricted to them, is a strong policy of depth at most $d-1$; granting the
additional projected observation cannot invalidate it. Its states are witnesses
for $\Gamma(q_{aj},S_{aj},d-1)$ because $T_j\subseteq S_{aj}$.
Thus $\card{T_j}\leq\CM(q_{aj},S_{aj},d-1)$ by induction. The actual assignments
of these distinct names to buckets form an integral feasible matching. Hence
$\card T\leq\card{F_a}+\nu_a\leq\max_{a'}(\card{F_{a'}}+\nu_{a'})$.
Also $\card T\leq\card S$. Maximize over the choices defining $\Gamma$ and
apply Lemma~\ref{lem:order}. This proof does not require a policy to succeed on
the entire child eligibility mask, nor does it identify a child capacity with
an arbitrary feasible subset of that cardinality.
\end{proof}

\begin{corollary}[Admissible bounded decisions]\label{cor:searchbound}
Let $S$ be the unfinished names of an exact belief with common $q$. A certified
upper bound $U$ on any capacity in Theorem~\ref{thm:sound} satisfying
$U<\card S$ proves $V(B)>d$. If the belief is already a goal its value is zero.
The first saturating depth of a complete recurrence is an admissible cost lower
bound; no saturation through $H$ proves only $V(B)>H$.
\end{corollary}
\begin{proof}
A strong policy for $B$ restricts to a successful policy on its unfinished
names. Therefore $V(B)\leq d$ would imply $\Gamma(q,S,d)=\card S$, contradicting
the certified upper bound. Already-goal names are removed only in this argument,
not from the actual search. The claim for the first saturating depth follows
by applying this implication at every smaller depth.
\end{proof}

\section{Hall deficiency: what sibling matching can certify}
Fix one action, put $R=R_a$, and let $\N(X)=\{j:E_{aj}\cap X\neq\varnothing\}$.
Use complete child matching capacities $c_j$ for the following exact identity.
\begin{theorem}[Deficiency characterization]\label{thm:hall}
\begin{align}
\nu_a&=\min_{X\subseteq R}\left(\card{R\setminus X}+
                                      \sum_{j\in\N(X)}c_j\right),\label{eq:cut}\\
\delta_a&=\card R-\nu_a
       =\max_{X\subseteq R}\left(\card X-\sum_{j\in\N(X)}c_j\right).\label{eq:deficit}
\end{align}
For a nonempty affordable action set,
$\CM(q,S,d)=\card S-\min_a\delta_a$.
Thus matching rejects exactly when every affordable action has a positive
Hall deficiency. If there is no affordable action, the capacity is zero.
\end{theorem}
\begin{proof}
Use source-to-name edges of capacity one, name-to-eligible-bucket edges of
capacity greater than $\card R$, and bucket-to-sink edges of capacity $c_j$.
The cut separating all names has capacity $\card R$, so a minimum cut does
not cut a name-to-bucket edge. If its source-side names are $X$, it must put
all adjacent buckets on the source side. Adding any other bucket cannot reduce
the cut. The resulting capacity is the expression in~\eqref{eq:cut}.
Conversely every $X$ defines such a cut. Integral max-flow/min-cut gives the
identity; subtract it from $\card R$ to get~\eqref{eq:deficit}. Since
$\card{F_a}+\card{R_a}=\card S$, the action maximum yields the last assertion.
The empty witness makes every deficiency maximum nonnegative.
\end{proof}

\begin{corollary}[Partial, locally checkable rejection]
Replace $c_j$ by any certified upper bounds $u_j\geq c_j$ and select any
$X_a\subseteq R_a$ for each affordable action. Then
\begin{equation}\label{eq:partialcut}
U=\min\left(\card S,\max_a\left[
 \card{F_a}+\card{R_a\setminus X_a}+
 \sum_{j\in\N(X_a)}u_j\right]\right)
\end{equation}
is an upper bound on $\CM$. If $U<\card S$, it is a sound rejection even if
some children remain unresolved and use the trivial $u_j=\card{S_{aj}}$.
\end{corollary}
\begin{proof}
Every feasible matching assigns at most $\card{R_a\setminus X_a}$ names
outside $X_a$ and at most the summed capacities of adjacent buckets from
inside $X_a$. Increasing capacities preserves this upper bound. Apply the
same argument to every affordable action, then take the maximum and cap.
\end{proof}
This is a use of a standard capacitated Hall/min-cut theorem, not a claim to
invent matching. A witness can be checked using intersections and sums in
$O(K\lceil N/w\rceil)$ word operations for $K$ buckets. A depth-decreasing DAG
of such bounds is checked in topological order, conditional on complete and
correct channel tables. Certifying those tables from symbolic formulas requires
separate SAT/enumeration evidence. No polynomial-size certificate in the compact
planning input is promised. The reference checker discovers cuts by exhaustive
enumeration only for tests; production matching already uses augmenting paths.

For example, buckets $\{0\},\{0\},\{1,2\}$ with capacities $1,1,1$ have
lineage sum three but matching value two. The witness $X=\{1,2\}$ has only
one adjacent bucket of capacity one. Proving that third capacity is enough;
the first two can retain their trivial unit upper bounds. This is a channel
example, not a newly solved public planning benchmark.

\section{Separation from complete subset and resource-information controls}
For an exact belief let $h_k(B)=\max_{\varnothing\neq T\subseteq B,\card T\leq k}
V(T)$, with the same resources and policy semantics. It is admissible because
a policy for $B$ restricts to every subset. This is an exact subset control,
not a claim to reproduce a published planner.

The second control constructs a graph with an edge between two worlds exactly
when no common legal action sequence of length at most $d$ reaches goals for
both under their resources. Observations are ignored in this synchronous test.
For each world $m$, let $K_m$ be the maximum number of observing actions in any
legal goal-reaching sequence of length at most $d$ with its resources; $K_m=-1$
if there is no such sequence. For every graph clique $Q$ require
\begin{equation}\label{eq:kraft}
\sum_{m\in Q}2^{-K_m}\leq1.
\end{equation}
This is the explicitly implemented resource-aware information comparator.

\begin{lemma}[Validity of the comparator]
Every feasible policy satisfies~\eqref{eq:kraft}, and has no $K_m=-1$.
\end{lemma}
\begin{proof}
Two endpoints of a conflict edge cannot share a terminal history: that history
would supply the prohibited common legal sequence. Clique members therefore
occupy distinct leaves. After contracting non-observing actions, their binary
observation words are prefix-free: a deterministic controller cannot both stop
and continue at the same history. If their actual observation depths are
$\ell_m$, Kraft's inequality gives $\sum 2^{-\ell_m}\leq1$. Their actual
goal-reaching paths satisfy $\ell_m\leq K_m$, which proves the inequality.
Single-child observations only weaken the test. Allowing already-goal states
to continue in the computation of $K_m$ avoids underestimating their paths.
\end{proof}

\begin{theorem}[Separation from the conjunction of complete controls]\label{thm:separate}
For every fixed $k\geq1$, there is an explicitly represented task with
$n>k$, $n\geq3$, candidates, one resource with budget two, and horizon two,
such that all the following hold:
\begin{enumerate}[label=(\roman*)]
\item Every nonempty proper candidate subset has a conformant plan of cost at
most two, so even testing all $2^n-2$ such subsets cannot reject this budget.
\item The complete resource-conflict graph is empty and every exact $K_m=1$.
Consequently every clique inequality passes.
\item On the same retained projection, $C_0=n$ and $\CL=\CM=n-1$.
The full task has no feasible strong policy with budget two.
\end{enumerate}
Thus the relation rejects a task passing the conjunction of the complete
$k$-subset, scalar, and defined resource-information controls.
\end{theorem}
\begin{proof}
Use hidden immutable identity $i\in\{0,\ldots,n-1\}$, initially true control
atom $start$, false atoms $p_0,\ldots,p_{n-1},end,done$, and goal $done$.
Retain $start,p_0,\ldots,p_{n-1},end$, but not $done$.
Every action costs one resource unit. Action $prepare_j$ requires $start$,
deletes it, sets $p_j$, and sets $done$ only for identity $j+1\pmod n$.
Action $finish_j$ requires $p_j$, deletes it, sets $end$, and sets $done$ for
every identity except $j$. The goal is otherwise preserved. Action $recover$
requires $end$ and sets $done$. Action $sense_i$ observes whether the identity
is $i$, has no physical effect, and is available at every stage.

For a proper subset choose an omitted $j$ and use $prepare_j;finish_j$.
For the full set a preparation followed by a finish leaves $j$ unfinished;
a preparation followed by sensing does not finish the remaining candidates.
Sensing first leaves a branch of at least two candidates and only one remaining
action, while one productive action at $start$ can finish only one identity.
Two senses achieve no goal; other first productive actions are inapplicable.
This exhausts all first actions and proves infeasibility, although an extra
unit permits $prepare_j;finish_j;recover$.

Every pair is a proper subset and has a common two-action goal sequence, so
the conflict graph is empty. A single world can sense and then choose its
successful preparation, giving $K_m\geq1$; it cannot achieve the goal with
two senses, giving $K_m=1$. Every nonempty clique is a singleton of weight
$1/2$, hence passes.

At $p_j$, the scalar immediate finish mask contains $n-1$ identities, including
the one already credited by preparation. Thus at the root scalar counting
gives $\min(n,1+(n-1))=n$. Lineage removes that first identity, so preparation
followed by finishing contributes $1+(n-2)=n-1$. At the root a sensing action
has two disjoint branches; with one remaining action each branch has capacity
one, so its contribution is two, at most $n-1$. These are all possible root
contributions. Hence $\CL=n-1$. Buckets in this construction are disjoint by
immutable identity, so matching has the same values. Taking
$n\geq\max\{3,k+1\}$ proves the claim for any fixed $k$.
\end{proof}

Given the explicit channel table for this family, only $O(n)$ relational
nodes at positive depths are reached: the root, the $n$ preparation successors,
and at most $2n$ sensing successor masks at $start$. Each scans $O(n)$ actions
with at most two buckets. Thus the combinatorial DP uses
$O(n^2\lceil n/w\rceil)$ word work, excluding construction and representation
of the channels. This does not bound general SAT construction time, and is not
an exponential runtime comparison against a well-implemented fixed-$k$ test.
Its point is the high-order infeasibility missed by \emph{every} proper subset.

\section{Projection refinement: a counterexample and qualified guarantees}
\begin{proposition}[More retained predicates can weaken the bound]\label{prop:counter}
There is a deterministic two-candidate task with no observations where both
relational capacities increase from one to two when one initially common
dynamic predicate is added to the projection, although the task is infeasible.
\end{proposition}
\begin{proof}
Let hidden immutable $h$ be true in one candidate and false in the other.
Initially $start$ is true and $stage,b,done$ are false. The only initial action
$prepare$ deletes $start$, sets $stage$, and sets $b$ exactly when $h$ is true.
There are two finishing actions, both requiring and deleting $stage$; one sets
$done$ when $h$ is true, the other when $h$ is false. There are no observations
or other actions. Each singleton has a two-step plan, but a policy cannot choose
different finishes for the indistinguishable candidates, and cannot finish twice.

For $P=\{start,stage\}$, preparation has one non-goal bucket containing both
names. The child has maximum immediate capacity one, so $\CL=\CM=1$ at depth
two. For $P'=P\cup\{b\}$, the initially common value is still false, but
preparation creates two different projected successors. The relaxation grants
their distinction, producing two singleton buckets each of capacity one.
Consequently both capacities are two. This is a counterexample in the actual
Boolean action language, not only in an arbitrary abstract table.
\end{proof}
In particular, neither relation universally dominates an exact two-world
control. In this refined example the relation passes, the singletons are
feasible, and the exact pair rejects. Together with
Theorem~\ref{thm:separate}, this rules out a universal dominance claim over
complete fixed-$k$ controls ($k\geq2$). The scalar/lineage/matching ordering
holds only when the projection and channels being compared are the same.

\begin{theorem}[A sufficient channel simulation]\label{thm:simulation}
Consider a fine and a coarse channel system on the same candidate names, and
a successor-compatible map $f$ from fine keys to coarse keys. Assume for each
fine key and action: (a) that action exists at its coarse image;
(b) the immediate goal masks are equal; and (c) there is an injective map
$\eta$ from its fine buckets to coarse buckets such that
$E'_j\subseteq E_{\eta(j)}$ and $f(q'_j)=q_{\eta(j)}$.
If the conditions hold recursively, then for $C=\CL$ or $\CM$,
$C'(q',S,d)\leq C(f(q'),S,d)$.
\end{theorem}
\begin{proof}
Induct on depth. Equal goal masks give equal immediate credit and the same
remaining mask. A fine child mask is a subset of its mapped coarse mask.
The induction hypothesis, followed by mask monotonicity, bounds its capacity
by the mapped coarse child. Injectivity lets us compare the lineage sums
without using a coarse contribution twice. For matching, send each fine
allocation to its mapped coarse bucket; its eligible edges and capacities are
available, and injectivity avoids merging two allocated child loads into one
capacity. Every fine action contribution is thus bounded by the corresponding
coarse one. Extra coarse actions, the maximum, and the equal mask caps preserve
the inequality.
\end{proof}
This is sufficient, not necessary, and deliberately stronger than syntactic
predicate inclusion. A refined observation split violates injectivity. Changing
existential goal masks also falls outside the theorem; it must not be silently
treated as satisfying the hypothesis. No production projection selector currently
checks this entire simulation condition.

\begin{corollary}[Safe envelope without a simulation assumption]
For the same exact belief and depth, take the minimum of any valid completed
capacity upper bounds from applicable projections. The result remains a sound
upper bound on the number of actual successful unfinished names and cannot
increase as more bounds are added. Equivalently, maximize their admissible
cost lower bounds. A stale certificate is reusable only for the same exact
belief/task/depth, not merely the same set of predicate identities.
\end{corollary}
\begin{proof}
Each applicable projection separately upper-bounds the actual successful
coverage by Theorem~\ref{thm:sound}. Their minimum therefore does too. Adding
another number to a minimum cannot increase it. The corresponding statement
about exact-belief lower-bound records follows directly.
\end{proof}

\begin{proposition}[What a multistep refinement core proves]
Let $F_\tau$ describe independently witnessed steps of a fixed candidate path.
For each forgotten predicate $p$, let $L_p$ equate its first value to the actual
initial value and each subsequent pre-value to the preceding post-effect value.
If $F_\tau\land\bigwedge_{p\in U}L_p$ is UNSAT, retaining and propagating
\emph{all} predicates in $U$ excludes that particular path pattern.
\end{proposition}
\begin{proof}
A refined path realizing the same action, observation, coarse-projection and
termination pattern supplies witnesses satisfying $F_\tau$. Propagation of all
retained $U$ supplies every $L_p$, contradicting UNSAT. An incomplete deletion
procedure may still produce a sufficient core, but not a minimum-cardinality
one. Choosing just one member of the core need not exclude the path.
\end{proof}
Excluding a path is not a proof that all relaxed paths disappear, that coverage
decreases, or that a real contingent policy is impossible. A fresh sound
capacity certificate is required before pruning.

\section{Partial computation and exact-search integration}
Let $v=(q,S,d)$ denote a DP node with interval $[L_v,U_v]$, initially
$[0,\card S]$ at positive depth and $[0,0]$ at depth zero or an empty mask.
An action with an unfinished channel table has interval $[0,\card S]$.
Once its full table is available, compute its lineage interval by substituting
child $L$ and $U$ into~\eqref{eq:lineage}; for matching use the two capacitated
networks, or any certified cut upper bound from~\eqref{eq:partialcut}.
Take the maximum of the action endpoints and apply the mask cap. Intersect
with the previous node interval, so $L$ only increases and $U$ only decreases.

\begin{theorem}[Safe interruption and decisive intervals]\label{thm:interval}
These updates enclose the complete recurrence under any sequence of pauses,
resumptions, and shared completed-cache hits. $U_v<\card S$ rejects the original
bounded task. $L_v=\card S$ proves only that this relaxation cannot reject it.
An action with upper endpoint at most the node's current lower endpoint can
be skipped when computing that node's exact maximum.
\end{theorem}
\begin{proof}
All action functions are monotone in child capacities. Replacing them by lower
and upper endpoints therefore encloses the true action value; maxima and
cardinality caps are also monotone. Unfinished tables retain the universal
interval, so unseen actions, goal memberships or outcomes cannot be omitted
from a rejection. Intersecting two valid intervals preserves validity.
Depth decreases on dependency edges, giving an inductive proof on the DAG
independent of evaluation order. A skipped action cannot exceed the already
established lower bound on the maximum. Rejection follows from
Corollary~\ref{cor:searchbound}; a saturated lower endpoint only establishes a
large optimistic capacity. A partial action maximum below $\card S$ is a lower
endpoint and has no rejection meaning.
\end{proof}
The current resumable lineage engine usually leaves $U_v=\card S$ until closing
the node. Its saved-child ceilings and dominated-action skipping are special
cases of this theorem. The stronger intermediate upper propagation and cut-DAG
rules above are proved and reference-tested, but are \emph{not} integrated into
the production scheduler. SAT UNKNOWN or incomplete outcome enumeration cannot
be treated as an empty bucket table.

\begin{theorem}[Sound and complete bounded search]
An exact depth-bounded AND/OR search remains sound and complete when it prunes
only with the above certified rejections or other admissible bounds, keeps
exact belief/resource keys, and uses only verified policies as upper bounds.
This holds under arbitrary query selection, provided optional computation is
bounded or fairly yields to search. Increasing thresholds finds the optimal
finite policy whenever the unrestricted search reaches its cost threshold.
\end{theorem}
\begin{proof}
Induct on remaining depth. Goal beliefs terminate correctly; a non-goal belief
at depth zero fails. A certified rejection excludes only infeasible bounded
subproblems. Otherwise enumerate actions legal in every world. An action
succeeds exactly when every nonempty actual observation child succeeds at one
less depth; induction gives soundness and completeness of this test.
If every action fails the bounded problem is infeasible. The action set and
depth are finite, so fair finite optional work cannot prevent termination.
Memoized lower records store $V(B)>d$ as a finite lower bound $d+1$, not global
infinity. Verified policies give genuine upper records. Neither needs heuristic
consistency for this depth-decreasing recurrence. An ancestor-targeted rejection
may unwind descendant frames, but may update only the certified belief; it
does not prove those descendants infeasible. Finally, exhaustive increasing
thresholds cannot skip an optimal solution using an admissible lower bound.
\end{proof}

\noindent A proof-oriented search interface is:
\begin{verbatim}
Decide(B, d):
    if all worlds in B satisfy goal: return a stop policy
    if d == 0: return NO
    if a verified policy for B has depth <= d: return it
    if a certified lower record for B is > d: return NO
    choose any applicable projection; optionally resume a finite query
    if its certified U < number of unfinished names:
        lower[B] = max(lower[B], d + 1); return NO
    for every action legal in every world of B:
        split exact successors only by actual observations
        solve every nonempty child at d - 1
        if all children have policies: return their joined policy
    lower[B] = max(lower[B], d + 1); return NO

Interrupting optional inference: retain its valid interval/cache,
then continue exact search. Interrupting exact search: return UNKNOWN.
\end{verbatim}
The pseudocode is bounded decision search, not a claim that scheduling by a
learned probability is optimal or that UNKNOWN is NO.

\section{Finite-resource keys and actual complexity}
\begin{proposition}[Exact finite-horizon resource normalization]
Let $c_i^{\max}=\max_a c_i(a)$. At remaining depth $d$, replacing
$r_i$ by $\min(r_i,d c_i^{\max})$ leaves all three complete capacities unchanged.
\end{proposition}
\begin{proof}
At depth zero all values are zero. For $d\geq1$, normalization preserves
affordability because $c_i(a)\leq c_i^{\max}\leq d c_i^{\max}$. For an
affordable action, clipping its normalized successor to the $(d-1)$ threshold
gives the same value as clipping its original successor. If $r_i\leq
d c_i^{\max}$ this is immediate. Otherwise both successors are at least
$(d-1)c_i^{\max}$. Boolean channel templates are independent of resources.
Induction on the normalized successor keys proves equality. The proof fails
if resource values affect goals, observations, effects, or replenishment.
\end{proof}

Let $A$ be the number of grounded actions, $p=\card P$, $K$ the maximum number
of non-goal buckets per action, $W=\lceil N/w\rceil$ bitset words, and $D$ the
number of relational nodes actually visited. Given complete cached channel
tables, lineage evaluation has conservative work
$O(D A K(W+b+1))$, including mask work and child-key/resource assembly.
Matching adds at most $O(D A N^2K)$ for the current simple augmenting-path
implementation. Channel construction, SAT, grounding, and formula manipulation
are additional costs; write them separately as $T_{\mathrm{channels}}$ and
$T_{\mathrm{front}}$. These expressions assume unit-cost bounded-size integer
and hash operations; bit complexity of large counters and dictionary behavior
can add further factors.

For an initial resource vector $r^0$ and depth limit $H$, a coarse bound on
the number of scalar keys is
\[
Q_H\leq\sum_{d=0}^{H}2^p\prod_{i=1}^{b}
             \bigl(1+\min\{r_i^0,d c_i^{\max}\}\bigr),
\qquad D\leq 2^N Q_H.
\]
Resource normalization can save repeated work but does not remove the
exponential candidate-mask factor. Boolean observations give
$K\leq2^{p+1}$. Sparse reachable masks and lazy channels may be much smaller;
these are empirical benefits, not polynomial guarantees for a lifted input.
Grounding can itself be exponential in schema arity, and a node budget alone
does not bound SAT construction time.

For a completed proof DAG with $E_D$ explicit child incidences, stored channel
data take $O(E_{\mathrm{channels}}W)$ words and DP entries/continuations take
$O(D(W+p/w+b)+E_D)$ words up to implementation overheads. A hot template hit
removes construction but not its $K$ bucket visits, intersections, resource
subtractions, and child-key lookups. Consequently counting only new cache
entries is not a bound on total computation. Time slices checked between SAT
calls can overrun by an atomic call; no hard real-time guarantee follows.

\section{Relation to prior work and exact claim boundaries}
Rintanen's $n$-distance framework~\cite{rintanen} establishes the importance of
small joint world sets. Its conformant distance must not be used unchanged as
an admissible contingent distance when sensing can help; the paper explicitly
discusses that restriction. Our comparator uses exact contingent subset costs.
Theorem~\ref{thm:separate} concerns that fully defined comparator, not all
algorithms in the earlier framework.

The LUG/POND work~\cite{bryce} already reasons with labels, overlapping support,
and cross-world interactions. Our proof does not establish dominance over its
full machinery. The candidate contribution is the particular resource-indexed
successful-name capacity and its certified integration into bounded search;
neither labels, max-flow, Kraft counting, nor choosing expensive heuristics is
claimed new by itself. A detailed equivalence/separation analysis against a
precisely instantiated labelled abstraction remains open. The present supplement
does not replace that comparison with a novelty assertion.

\section{Proof-to-code map and verification status}
\begin{description}[style=nextline]
\item[Implemented before this supplement.]
\texttt{symbolic\_coverage.py} constructs existential channels and $C_0$;
\texttt{relational\_coverage.py} implements both relations;
\texttt{bounded\_resource\_coverage.py} normalizes resources;
\texttt{resumable\_capacity.py} and \texttt{catalog\_capacity.py} resume lineage
computation; \texttt{threshold\_search.py} uses exact beliefs and finite lower
records; \texttt{resource\_information.py} implements the defined comparator;
\texttt{coverage\_separation.py} generates the separation family.
\item[New reference checks, not a production algorithm change.]
\texttt{coverage\_certificates.py} checks local lineage/Hall upper steps assuming
certified children and complete channels. \texttt{test\_theory\_strengthening.py}
checks the cut formula against independent assignments, the stronger separation,
mask/depth monotonicity, and the concrete refinement counterexample. It tests
8,539 matching networks and 66,102 cut bounds; all 238 nonempty proper subsets
of the $n=3,\ldots,7$ families; 55 completed resource-aware pair checks; and
3,456 mask-monotonicity comparisons. These finite checks can find mistakes but
are not machine-checked mathematical proofs or held-out performance evidence.
\item[Proved here but not integrated as a new production mechanism.]
General intermediate upper propagation, Hall-cut proof DAGs, and the sufficient
channel-simulation checker. Exact-belief maximum lower records already preserve
previous certificates; the implementation does not enumerate a portfolio of
all projections. No runtime improvement is inferred from these proofs.
\item[Still unproved or unsupported.]
Universal dominance over labelled planning heuristics or resource-aware bounds;
monotone strengthening from arbitrary added predicates; polynomial total
complexity; calibrated scheduling probabilities; net runtime gains on held-out
benchmarks; guarantees for stochastic dynamics, replenishment, numeric goals,
or an inaccurate LLM-supplied action model. None is needed as an assumption
hidden inside the theorems above.
\end{description}

\begin{thebibliography}{9}
\bibitem{rintanen} Jussi Rintanen. Distance Estimates for Planning in the Discrete
Belief Space. AAAI, 2004.
\url{https://users.aalto.fi/~rintanj1/jussi/papers/Rintanen04heu.pdf}.
\bibitem{bryce} Daniel Bryce, Subbarao Kambhampati, and David E. Smith.
Planning Graph Heuristics for Belief Space Search. Journal of Artificial
Intelligence Research 26:35--99, 2006.
\url{https://doi.org/10.1613/jair.1869}.
\end{thebibliography}
\end{document}
